\section{Preliminaries}
\label{sec:problem}

\smallskip\noindent\textbf{Problem definition.}
We are given a relational dataset $\mathcal{D}$ and an associated \emph{normative document}. This textual source could be laws, regulations, policy manuals, or operational specifications. It defines the requirements the data must satisfy, expressed in natural language rather than formal, machine-executable constraints.

We model a normative document as a set $\mathcal{S}$ of \emph{segments}, with each segment $s \in \mathcal{S}$ representing a distinct, self-contained textual unit in the document. The segmentation granularity may vary. A segment can correspond, for example, to a single sentence, multiple sentences, or a larger passage. We consider a document normative if it contains at least one segment that expresses a requirement whose compliance can be assessed based on the data.





%We are given a relational dataset $\mathcal{D}$ and an associated \emph{normative document} $\mathcal{N}$, a textual source that defines rules or requirements that data must satisfy within a given domain. Such documents include laws, regulations, policy manuals, and operational specifications, and typically express such requirements in natural language rather than as formal, machine-executable constraints.

%A provision is a textual unit within a normative document that may express a checkable  requirement governing the data. 

%$s \in \mathcal{T}$ can range from a single sentence to a group of sentences, a paragraph, a %section, an entire document, or any other desired text unit.

A \emph{data quality constraint} is a formal rule that $\mathcal{D}$ must satisfy (we discuss the choice of constraint formalism below).
Constraints are evaluated on a target relation $r \in \mathcal{D}$ with schema $R = (A_1, \ldots, A_n)$, where each $A_i$ is an attribute with domain
$\mathrm{dom}(A_i)$. If a constraint involves multiple relations, we evaluate it on a derived relation obtained by applying the required joins.


Let $\Phi$ denote the space of constraints over $R$.
The \emph{Text-to-Constraints} problem is to extract a set of segment--constraint pairs
\[
\Pi = \{(s_1, \varphi_1), \ldots, (s_k, \varphi_k)\} \subseteq \mathcal{S} \times \Phi,
\]
where $s_i \in \mathcal{S}$ is a segment and $\varphi_i \in \Phi$ is a constraint whose semantics are grounded in the normative content of $s_i$.
Each pair must satisfy two requirements: \emph{executability}, ${\varphi_i}$ can be evaluated on any instance of $R$, 
%(the constraint can be evaluated on an instance of  $\mathcal{T}$) 
and \emph{traceability}, $s_i$ is identifiable as the textual justification for $\varphi_i$.
%(the system can identify a textual justification for $\varphi$).

%In this work, we consider two granularities: sentence-level extraction, where each provision is a single sentence and yields at most one constraint, and document-level extraction, where the entire document is treated as one provision and may yield multiple constraints. More generally, one provision may yield multiple constraints and a constraint may depend on multiple provisions; such settings and other granularities are left for future work.

In general, one segment may yield multiple constraints, and a constraint may depend on multiple segments. In this work, we consider \DIFdelbegin \DIFdel{two granularities: }\DIFdelend sentence-level extraction, where each segment is a single sentence and yields at most one constraint\DIFdelbegin \DIFdel{, and document-level extraction, where }\DIFdelend \DIFaddbegin \DIFadd{; coarser granularities, in which }\DIFaddend constraints are derived by jointly reasoning over \DIFdelbegin \DIFdel{multiple segmentswithin a normative document}\DIFdelend \DIFaddbegin \DIFadd{several segments, are left for future work}\DIFaddend .


\smallskip\noindent\textbf{Denial constraints.}
We instantiate $\Phi$ as the space of \emph{denial constraints} (DCs). A DC specifies that a conjunction of predicates must never be simultaneously true. Formally, given a relation instance $r$ over schema $R$, a DC $\varphi$ is a first-order sentence of the form:
\[
\varphi:\ \forall\, t_1,\ldots,t_m \in r:\ \neg\big(p_1 \land \cdots \land p_\ell\big),
\]
where each $p_i$ is an atomic predicate over tuple variables and constants, and $m \ge 1$ is the number of tuple variables. While DCs may in principle involve an arbitrary number of tuple variables, we focus on DCs with $m \le 2$. This restriction is common in the literature and adopted by prior work~\cite{chu2013discovering,martin2025false, Pena2019DiscoveryDC}, as DCs involving more than two tuple variables are rare in practice and lead to a 
dramatic increase in both search and semantic complexity.


Let $A, B \in R$ and let $c$ denote a constant in $\mathrm{dom}(A)$. We allow the following predicate forms: (i) comparisons $u\,\theta\,v$ with $\theta \in \{=,\neq,<,\leq,>,\geq\}$ where $u$ and $v$ are tuple attributes (e.g., $t.A$, $t'.A$, $t.B$) or constants; (ii) pattern predicates $t.A\ \mathsf{LIKE}\ p$ and $t.A\ \mathsf{NOT\ LIKE}\ p$ where $p$ is an SQL pattern; (iii) set-membership predicates $t.A\ \mathsf{IN}\ C$ and $t.A\ \mathsf{NOT\ IN}\ C$ where $C \subseteq \mathrm{dom}(A)$; and
(iv) null checks $t.A\ \mathsf{IS\ NULL}$ and $t.A\ \mathsf{IS\ NOT\ NULL}$.

A DC $\varphi$ is \emph{violated} if there exist tuples for the quantified variables such that all predicates hold simultaneously.

%When a DC uses two-tuple variables, it corresponds to a self-join query. For example, 
DCs can be translated directly to SQL.
For example,  the SQL query in Example~\ref{ex:intro} is actually equivalent to the following DC:

\begingroup\small
\[
\varphi:\ \forall\, t, t' \in \mathit{applications}:
\]
\[ 
\neg\!\left(
\begin{aligned}
& t.\mathit{year}=t'.\mathit{year}\ \land\ t.\mathit{edition}=t'.\mathit{edition}\ \land \\
& t.\mathit{institution\_id}=t'.\mathit{institution\_id}\ \land\ t.\mathit{campus\_id}=t'.\mathit{campus\_id}\ \land \\
& t.\mathit{course\_code}=t'.\mathit{course\_code}\ \land\ t.\mathit{shift}=t'.\mathit{shift}\ \land \\
& t.\mathit{application\_mode}=t'.\mathit{application\_mode}\ \land \\
& t.\mathit{score}>t'.\mathit{score}\ \land\ t.\mathit{ranking}>t'.\mathit{ranking}
\end{aligned}
\right).
\]
\endgroup

\smallskip\noindent\textbf{Why DCs.}
DCs provide a compact formalism for a broad class of integrity constraints, including checks, keys, functional dependencies, order dependencies, conditional variants, and other cross-row conditions. They are also directly checkable on a relation through their violation semantics, making them a natural target for compliance-oriented rule extraction. Compared to SQL, which allows many syntactically different but equivalent formulations of the same rule, DCs offer a more uniform and canonical representation, yielding a more stable target space for extraction and evaluation. Finally, with the standard restriction to at most two tuple variables, the space of candidate constraints remains manageable in practice. Future work can focus on extending the proposed framework to even richer formalisms and/or SQL encodings.
% DCs offer a compact formalism for a broad class of integrity constraints (e.g., checks, keys, functional dependencies, ordering constraints, conditional variants, and cross-row
% conditions) and can be checked directly on a relation.
%Compared to prior DC-based work that typically restricts predicates to basic comparisons~\cite{chu2013discovering}, we also support patterns (\texttt{LIKE}/\texttt{NOT
%LIKE}), membership (\texttt{IN}/\texttt{NOT IN}), and null checks, which are often required to encode normative text. To keep the search space under control, we adopt the standard restriction of at most two tuple variables.
% SQL, by contrast, admits many different syntactic formulations of the same rule (e.g., using self-joins, \texttt{GROUP BY}/\texttt{HAVING}, or correlated subqueries),
% requiring exploration of numerous alternative rewritings of the same constraint. DCs provide a uniform representation instead.
% Future work can extend the framework to richer predicates or SQL encodings while retaining a similar workflow.







%%% old stuff
%We are given a relational dataset $\mathcal{D}$ and its associated text $\mathcal{T}$
%(e.g., regulations, policies, manuals).
%We model $\mathcal{T}$ as a set of text segments, where a segment $s \in \mathcal{T}$ can range from a single sentence to a group of sentences, a paragraph, a section, an entire document, or any other desired text unit.
%
%A \emph{data quality constraint} is a formal rule that the data must satisfy (we discuss
%the choice of constraint formalism below).
%Constraints are evaluated on a target relation $r$ with schema $R = (A_1, \ldots, A_n)$,
%where each $A_i$ is an attribute with domain $\mathrm{dom}(A_i)$. If a constraint involves
%multiple relations, we evaluate it on a derived relation obtained by applying the required
%joins.